\section{Actual normalized spectral geometry and leave-one-out control}
\label{sec:geometry}

This section proves the spectral premises used by both fitted decoders.
Leave-one-out (LOO) matrices occur only in the proof. The implemented
encoder retains one full-graph eigensystem.

\subsection{Raw concentration and deterministic spectral consequences}

For any prescribed fixed $H>0$, when $d\ge c\log n$ there is a fixed
$L_H$ such that
\begin{equation}\label{eq:raw-norm-tail}
 \PP\{\norm{A-\Omega}_{\op}>L_H\sqrt d\}\le e^{-Hd}.
\end{equation}
Indeed, apply Corollary 3.12 and its centered-entry extension in
Remark 3.13 of Bandeira--van Handel~\cite{bvh} to $A-\E A$.
For bounded centered entries that result bounds the tail at
$C\sqrt{\max_i\sum_j\operatorname{Var}(A_{ij})}+t$ by
$n\exp(-t^2/C')$. Here the row variance sum is $O(d)$, and the entries
have absolute value at most one. Choosing $t$ a sufficiently large fixed
multiple of $\sqrt d$ absorbs both $H$ and $\log n$.
The completed-mean diagonal costs $O(p)$. The centered-entry extension is
needed: centered Bernoulli noise does not have a symmetric distribution.

Fix constants $a,b,L>0$ and define the base event
\begin{equation}\label{eq:base-event}
 \mathcal E_0=
 \{\norm{A-\Omega}_{\op}\le L\sqrt d,
       \quad ad\le D_i\le bd\text{ for every }i\}.
\end{equation}
Its minimum-degree probability is deliberately not included in
\eqref{eq:raw-norm-tail}; it is treated separately in
\cref{sec:probability}.
On $\mathcal E_0$,
\begin{equation}\label{eq:degree-aggregate}
 \norm{(D-\bar D)\one}_2\le L\sqrt{nd},\qquad
 \sum_j\frac{(D_j-\bar D_j)^2}{D_j}\le Cn,
 \quad \bar D=\diag(\Omega\one).
\end{equation}
These follow by applying $A-\Omega$ to $\one$; no uniform vanishing
relative degree error is required.

The observed-degree-normalized completed mean
\[
 M=D^{-1/2}\Omega D^{-1/2}
\]
has rank $K$, column space $\col(D^{-1/2}\Theta Z)$, and its nonzero
eigenvalues have absolute values at least a fixed positive constant.
For the last statement,
$Z^\top\Theta D^{-1}\Theta Z$ is diagonal with entries comparable to
$n/d$, and the nonzero eigenvalues are those of
\[
 p(Z^\top\Theta D^{-1}\Theta Z)^{1/2}
 C(Z^\top\Theta D^{-1}\Theta Z)^{1/2}.
\]
The lower singular-value assumption supplies the fixed gap.
Since $\norm{S-M}_{\op}\le C/\sqrt d$, the retained eigenvalues are
bounded away from zero, the discarded spectrum is $O(d^{-1/2})$, and
\begin{equation}\label{eq:normalized-angle}
 \norm{UU^\top-\Pi_M}_{\op}\le C/\sqrt d,
\end{equation}
where $\Pi_M$ projects onto $\col(D^{-1/2}\Theta Z)$.
These gap estimates follow directly by expanding a retained vector in the
signal and complementary eigenbases, or by the usual subspace perturbation
inequality. Positive and negative signal eigenvalues are both retained.

The same conclusions hold for every vertex-deleted graph with its degrees
\emph{recomputed}: deleting a vertex preserves the raw principal-submatrix
norm bound, reduces each surviving degree by at most one, and changes a
block mass by only a bounded amount. All constants can be chosen common
for the full graph and these $n$ minors.

\subsection{Exact generalized equations for one incident row}

Put the tested vertex first and write
\begin{equation}\label{eq:arrowhead-A}
 A=\begin{pmatrix}0&x^\top\\x&B\end{pmatrix},\qquad
 k=\norm{x}_1,\quad \Delta=\diag(B\one),\quad
 D=\diag(k,\Delta+E_x),\quad E_x=\diag(x).
\end{equation}
Here $x\in\{0,1\}^{n-1}$. Suppose the deleted degrees lie between
fixed multiples of $d$, $0<k\le bd$, and
\[
 S_B=\Delta^{-1/2}B\Delta^{-1/2}
    =U_B\Lambda_BU_B^\top+N,
 \quad U_B^\top U_B=I_K,\quad NU_B=0,
\]
with signal eigenvalues separated from zero and
$\norm{N}_{\op}\le C/\sqrt d$. Write
\[
 V=\Delta^{-1/2}U_B,\quad Q_\perp=I-U_BU_B^\top,\quad
 F=\Delta^{-1/2}E_x\Delta^{-1/2},\quad h=\Delta^{-1/2}x.
\]
Then $\norm F\le C/d$ and $\norm h^2\le Ck/d$.
The full normalized matrix differs from $\diag(0,S_B)$ by
$O(d^{-1/2})$: its off-diagonal normalized row has squared norm at most
$C/d$, and the lower degree-renormalization perturbation is $O(d^{-1})$.
The full retained signal therefore has the same fixed gap.

Take full generalized eigenvectors $g_r=D^{-1/2}u_r$, so
$Ag_r=\mu_rDg_r$ and jointly $G^\top DG=I_K$.
Decompose them exactly as
\[
 g_r=\binom{t_r}{\Delta^{-1/2}(U_By_r+z_r)},
 \qquad U_B^\top z_r=0.
\]
Let $h_V=U_B^\top h$ and $h_Q=Q_\perp h$.
The top row and complementary equations give
\begin{align}
 h^\top(U_By_r+z_r)&=\mu_r k t_r,\label{eq:generalized-top}\\
 (\mu_rI-N+\mu_rQ_\perp FQ_\perp)z_r
   &=t_rh_Q-\mu_rQ_\perp FU_By_r.\label{eq:generalized-complement}
\end{align}
The inverse $T_r$ of the complementary operator has bounded norm.
Substituting into \eqref{eq:generalized-top} gives
\begin{equation}\label{eq:top-resolvent}
 (\mu_r k-h_Q^\top T_rh_Q)t_r
 =h_V^\top y_r-\mu_rh_Q^\top T_rQ_\perp FU_By_r.
\end{equation}
Its denominator has magnitude at least $ck$: the second term is at most
$Ck/d$. This remains true when $k$ is small.

Put $M_B=\norm{U_B}_{2\to\infty}$.
The coordinates $y_r$ have bounded norm because
$\Delta^{1/2}G_{-i}=(I+F)^{-1/2}U_{-i}$ has norm at most one.
Using
\[
 \norm{h_V}\le CkM_B/\sqrt d,\qquad
 \norm{FU_B}_F\le CM_B\sqrt k/d
\]
in the exact equations yields
\begin{equation}\label{eq:generalized-bounds}
 |t_r|\le CM_B/\sqrt d,\qquad
 \norm{z_r}_2\le CM_B\sqrt k/d.
\end{equation}
The matrix $Y=(y_1,\ldots,y_K)$ has bounded inverse: the full and padded
deleted signal spaces have angle $O(d^{-1/2})$, and the extra multiplier
$(I+F)^{-1/2}$ costs only $O(d^{-1})$.

\begin{lemma}[Normalized row bounds from model events]\label{lem:normalized-rows}
Under the structural model assumptions, for any prescribed $H>0$ there
are row-concentration events of total failure probability at most
$2Kn e^{-Hd}$ such that, on their intersection with $\mathcal E_0$,
the full normalized retained vectors and every deleted normalized retained
basis satisfy
\begin{equation}\label{eq:row-incoherence}
 \norm{U}_{2\to\infty}\le C_H/\sqrt n,
 \qquad \max_i\norm{U^{(-i)}}_{2\to\infty}\le C_H/\sqrt n.
\end{equation}
They also satisfy the gap and angle bounds above. In particular, with
$\Pi_0=Z(Z^\top Z)^{-1}Z^\top$,
\begin{equation}\label{eq:euclidean-angle}
 \norm{\Pi_X-\Pi_0}_{\op}\le C_H/\sqrt d,
 \quad \norm{Q}_{2\to\infty}\le C_H/\sqrt n,
\end{equation}
and the same holds for all recomputed deleted projectors.
\end{lemma}

\begin{proof}
Let $M_\mathrm{full}=\norm U_{2\to\infty}$ and
$M_i=\norm{U^{(-i)}}_{2\to\infty}$.
The lower rows of the full eigenbasis are
$(I+F)^{1/2}(U_BY+Z_\perp)$, with
$Z_\perp=(z_1,\ldots,z_K)$.
Using $Y^{-1}$ and \eqref{eq:generalized-bounds} gives
\[
 M_i\le C M_\mathrm{full}+CM_i/\sqrt d,
 \qquad M_i\le C' M_\mathrm{full}\quad\text{for large }d.
\]
This argument has not assumed row incoherence.
For the tested full row, \eqref{eq:top-resolvent}, $k\asymp d$ and
$u_{i,r}=\sqrt{k}\,t_r$ imply the sharper estimate
\begin{equation}\label{eq:tested-row-bootstrap}
 \norm{U_{i\cdot}}_2
 \le\frac C{\sqrt d}\norm{x^\top\Delta^{-1/2}U_B}_2+CM_i/d.
\end{equation}

Conditional on $B$, its eigenbasis and recomputed degrees are fixed and
$x$ is an independent Bernoulli row. Its mean action has norm at most
$C\sqrt{d/n}$, by $\norm{\E x}_2\le Cp\sqrt n$ and
$\norm{\Delta^{-1/2}U_B}\le C/\sqrt d$.
We use the following weighted Bernoulli bound from
Abbe--Fan--Wang--Zhong~\cite[Lemma 7, Eq.\ (C.2)]{afwz}:
for fixed $v$ and independent Bernoulli coordinates with success
probabilities at most $q$,
\begin{equation}\label{eq:afwz-row}
 \PP\left\{\left|\sum_jv_j(X_j-\E X_j)\right|
 >\frac{(2+\alpha)nq\norm v_\infty}
 {1\vee\log(\sqrt n\norm v_\infty/\norm v_2)}\right\}
 \le2e^{-\alpha nq}.
\end{equation}
Apply it to each of the fixed $K$ columns of $\Delta^{-1/2}U_B$, taking
$q=\min(Cp,1)$, which is comparable to $p$ on the structural class,
and selecting $\alpha$ to give rate $Hd$.
The column norms are at most $C/\sqrt d$ and $CM_i/\sqrt d$ in
$\ell_2$ and $\ell_\infty$, respectively. The function
$b/(1\vee\log(\sqrt n b/s))$ is nondecreasing in both admissible norm
arguments. Thus outside conditional probability $2Ke^{-Hd}$,
\[
 \norm{x^\top\Delta^{-1/2}U_B}_2
 \le C\sqrt{d/n}+
 \frac{C_H\sqrt d\,M_i}{1\vee\log(c\sqrt nM_i)}.
\]
A union over $i$, followed by \eqref{eq:tested-row-bootstrap}, yields
\[
 M_\mathrm{full}\le C/\sqrt n+
 \frac{C_H'M_\mathrm{full}}
 {1\vee\log(c'\sqrt nM_\mathrm{full})}
 +C'M_\mathrm{full}/d.
\]
For large $d$, absorb the last term. If $\sqrt nM_\mathrm{full}$ were
larger than a sufficiently large fixed constant, the logarithmic term
could also be absorbed, giving a contradiction. This proves
\eqref{eq:row-incoherence} and its deleted versions.

The conditioning is legitimate even though $\mathcal E_0$ depends on
$x$: define the conditional row event for every deleted graph satisfying
its measurable degree and gap premises, apply \eqref{eq:afwz-row} there,
and use $\mathcal E_0$ only for deterministic implications afterwards.

For \eqref{eq:euclidean-angle}, bounded-condition-number multiplication by
$D^{-1/2}$ in \eqref{eq:normalized-angle} compares $\col(X)$ with
$\col(D^{-1}\Theta Z)$ at angle $O(d^{-1/2})$.
Multiplying block $a$ of $D^{-1}\Theta Z$ by $dS_a/\sqrt{n_a}$ gives
entries $(\bar D_i/D_i)/\sqrt{n_a}$ on that block. By
\eqref{eq:degree-aggregate}, this basis differs from
$Z\diag(n_a^{-1/2})$ in Frobenius norm by $O(d^{-1/2})$.
This proves the angle to the block space. Finally,
$Q=X(X^\top X)^{-1/2}$ has the stated row bound since the eigenvalues of
$X^\top X=U^\top D^{-1}U$ are comparable to $1/d$.
The argument applies to DC-SBM as well as SBM.
\end{proof}

\subsection{A deterministic full-to-LOO score lemma}

\begin{lemma}[Both actual full-graph reconstructions have bounded LOO error]
\label{lem:loo}
Assume the deleted degree, spectral gap and discarded-spectrum conditions
following \eqref{eq:arrowhead-A}, as well as
$\norm{U_B}_{2\to\infty}\le C/\sqrt n$ and $k\le bd$.
For every $w\in\R^n$ with $\norm w_\infty\le W$, let
\[
 \Pi_B=V(V^\top V)^{-1}V^\top,
 \qquad H_B=VV^\top\Delta,\qquad V=\Delta^{-1/2}U_B.
\]
Uniformly over such incident rows and coefficient vectors,
\begin{align}
 (R_Ew)_i&=x^\top\Pi_Bw_{-i}+O(1),\label{eq:loo-euclidean}\\
 (R_Dw)_i&=x^\top H_Bw_{-i}+O(1).\label{eq:loo-weighted}
\end{align}
The constants depend only on the displayed fixed bounds and $W$.
For the hollow degree-metric reconstruction on the model good event,
\eqref{eq:loo-weighted} also holds with $R^\circ$ in place of $R_D$.
The right-side coefficients use degrees recomputed from $B$. When $w$
is fixed, as in the oracle contrasts, or is measurable from $B$, those
coefficients are independent of the tested row. Uniform scalar stability
alone does not make a full-row-dependent $w$ independent.
\end{lemma}

\begin{proof}
The case $k=0$ has zero row action and can be given that convention.
For $k>0$, \eqref{eq:generalized-bounds} now yields
\begin{equation}\label{eq:loo-refined}
 \norm t\le\frac C{\sqrt{nd}},\qquad
 \norm{\Delta^{-1/2}Z_\perp}_{\op}
 \le\frac{C\sqrt k}{d^{3/2}\sqrt n}.
\end{equation}
For the Euclidean action right-multiply $G$ by
$Y^{-1}(V^\top V)^{-1/2}$. This does not change its column space and
gives a basis
\[
 \mathcal F=\binom{\alpha}{Q_B+E},\quad
 Q_B=V(V^\top V)^{-1/2},\quad
 \norm\alpha\le C/\sqrt n,\quad
 \norm E\le C\sqrt k/(d\sqrt n).
\]
Its Gram matrix is $I+J$, where
$\norm J\le C\{\sqrt k/(d\sqrt n)+1/n\}$.
The deleted Euclidean basis has row norm at most $C/\sqrt n$, so
\begin{align*}
 \norm{x^\top Q_B}&\le Ck/\sqrt n,&
 \norm{x^\top E}&\le Ck/(d\sqrt n),\\
 \norm{Q_B^\top w_{-i}}&\le W\sqrt n,&
 \norm{\mathcal F^\top w-Q_B^\top w_{-i}}
 &\le C_W(1/\sqrt n+\sqrt k/d).
\end{align*}
Expand
\[
 (R_Ew)_i=(x^\top Q_B+x^\top E)(I+J)^{-1}\mathcal F^\top w
\]
about $x^\top Q_BQ_B^\top w_{-i}$.
The errors are bounded by a constant times
\begin{equation}\label{eq:loo-errors}
 k/d+k/n+k\sqrt k/(d\sqrt n)=O(1),
\end{equation}
because $k\le bd$ and $d\le Cn$. This proves
\eqref{eq:loo-euclidean}. A coarse $O(n^{-1/2})$ bound for the lower
rows alone would not suffice; the stronger $E$ bound is essential.

For the degree-metric action, put $v_r=\Delta^{-1/2}z_r$,
$f_r=Vy_r+v_r$ and $w_V=V^\top\Delta w_{-i}$.
The exact top equation gives $k\mu_rt_r=x^\top f_r$, whence
\begin{align}\label{eq:weighted-row-expansion}
 (R_Dw)_i
 =\sum_r(x^\top Vy_r+x^\top v_r)
 \{kt_rw_i+y_r^\top w_V+v_r^\top\Delta w_{-i}
                              +f_r^\top E_xw_{-i}\}.
\end{align}
The normalization $G^\top DG=I$ gives
\[
 Y^\top Y+Z_\perp^\top Z_\perp+ktt^\top
 +(VY+\Delta^{-1/2}Z_\perp)^\top E_x
            (VY+\Delta^{-1/2}Z_\perp)=I.
\]
All terms except $Y^\top Y$ have norm at most $Ck/(nd)$, by
\eqref{eq:loo-refined} and the row bound on $V$.
Since $Y$ is square, $\norm{YY^\top-I}\le Ck/(nd)$ as well.
The leading term in \eqref{eq:weighted-row-expansion} is
$x^\top VYY^\top w_V$; its difference from
$x^\top VV^\top\Delta w_{-i}$ is at most $C_Wk^2/(nd)$.
The other factors satisfy
\begin{align*}
 \norm{x^\top V}&\le Ck/\sqrt{nd},&
 \norm{w_V}&\le C_W\sqrt{nd},\\
 \norm{(x^\top v_r)_r}&\le Ck/(d^{3/2}\sqrt n),&
 \norm{(v_r^\top\Delta w_{-i})_r}&\le C_W\sqrt{k/d},\\
 \norm{(f_r^\top E_xw_{-i})_r}&\le C_Wk/\sqrt{nd},&
 \norm{kt}&\le Ck/\sqrt{nd}.
\end{align*}
Their products are bounded by constants times
$k/d$, $k^2/(nd)$ and $k\sqrt k/(d\sqrt n)$, all bounded.
This proves \eqref{eq:loo-weighted} with negative as well as positive
informative eigenvalues.

On the model good event, full row incoherence implies
$|(H_D)_{ij}|\le C/n$, $|(R_D)_{ii}|\le Cd/n=Cp$ and
$|s_i|\le Cp$. The hollow adjustment to a bounded row contrast is
$s_i\{(H_Dw)_i-w_i\}=O(p)$, proving the last assertion.
\end{proof}

The result compares the \emph{actual} full-graph normalized algorithm to
a row-independent score, with its changing degree normalizers included.
It does not treat a full-graph eigenvector as independent of the row that
created it. This is the dependence step needed before any conditional
sharp-tail argument is applied.
